\section{Changing the input shift without adding coordinates}
\label{sec:p3-beta-shift}

The substitution $\beta=X^2+2$ is convenient over $\mathbb Q$, but its
interaction with the five scalar contacts is not stable under extension
of the ground field. We first describe a uniform rational modification
and then a field-dependent way to avoid every member of a prescribed
finite exceptional set. Neither statement changes the elliptic
cyclic-group hypothesis.

\begin{proposition}\label{p3:integer-shift}
Let $C$ be an integer with $C\ge2$ such that $2401-C$ is not a
nonnegative integer square. Replace $\beta=X^2+2$ by
$\beta=X^2+C$ throughout the shared-point construction and the direct
Lucas-polynomial refinement. Under their stated arithmetic inputs,
the resulting rational input projection is still $\mathbb Z$.
The retained-coordinate counts and finite-algebra ranks are unchanged.
In particular every integer $C\ge2402$ is permitted.
\end{proposition}
\begin{proof}
Write $X=a/b$ in lowest terms, with $b>0$. Then
\[
                \beta=\frac{a^2+Cb^2}{b^2},
          \qquad \gcd(a^2+Cb^2,b^2)=1.
\]
Thus $\beta$ is an integer if and only if $X$ is an integer.
Moreover, at any finite prime where $\beta$ has negative valuation,
$v_p(X)<0$ and $v_p(\beta)=2v_p(X)$ is even. These are precisely
the input-specific facts used in scalar soundness: the finite
five-contact radical accounts for every odd contact, without a new
contact at infinity. The scalar height argument and the comparisons
which force $\beta$ to equal an integer label therefore remain valid.

For integral $X$, the new value of $\beta$ is an integer at least
two. It cannot equal $0,1,-343/32,$ or $1/2$. It cannot equal
$2401$, by the hypothesis on $C$. Hence the scalar contact product
\[
 r_\beta=\beta(\beta-1)(\beta-2401)(32\beta+343)(2\beta-1)
\]
is nonzero. The scalar normalization, local squareclass choices,
prime-pattern construction, and simultaneous point choices apply
to this integer value of $\beta$ just as before. Their estimates
were uniform for all integral $\beta\ge2$ away from these contacts;
they did not require the special numerical value of the constant
two. In particular the same height constant and precision choices
are valid. This proves completeness.

Changing a fixed coefficient introduces no coordinate and no new
algebra generator. It also does not change the formal degrees of the
source expressions. The independently computed upper bounds for the
emitted coefficient formulas must nevertheless be recorded separately,
since their numerical coefficient sizes can change.
\end{proof}

If a number field $K$ has a real embedding and $C\ge2402$, then
$X^2+C$ avoids every scalar contact for every $X\in K$. Indeed at
that embedding its value is at least $2402$, whereas every contact
is real and at most $2401$. The normalizing denominator
$\beta^2+\beta+23$ is also nonzero: its real discriminant is
$-91$. Thus $C=2402$ removes these two specific domain obstructions
over every number field with a real embedding. This observation
does not imply that the resulting formula defines $\mathcal O_K$:
the cyclic point group and the arithmetic comparison arguments
require their own field-dependent replacements.

\begin{lemma}[Avoiding a finite set by a translated square]
\label{p3:finite-square-avoidance}
Let $K$ be a number field and let $\Theta\subset K$ be finite.
There is an effectively searchable arithmetic progression of positive
integers $C$ such that
\[
                  \theta-C\notin K^2
                       \qquad(\theta\in\Theta).
\]
In particular $X^2+C$ avoids $\Theta$ for every $X\in K$. The
progression contains integers greater than any prescribed bound.
\end{lemma}
\begin{proof}
For each $\theta\in\Theta$, choose a different odd rational prime
$\ell_\theta$ which splits completely in the normal closure of
$K/\mathbb Q$ and at which $\theta$ is integral. Infinitely many
such primes exist by Chebotarev's theorem
\cite[Section~2.1, Theorem~1]{p2lbSerre1981}. Excluding the finitely
many ramified primes, denominators, and previously used primes
does not affect existence. Choose a prime $\mathfrak p_\theta$
of $K$ above $\ell_\theta$. Its residue field is
$\mathbb F_{\ell_\theta}$.

Choose a nonzero nonsquare
$u_\theta\in\mathbb F_{\ell_\theta}$ and impose the integer
congruence
\[
                C\equiv\overline\theta-u_\theta
                       \pmod{\ell_\theta}.
\]
Every such $C$ makes $\theta-C$ a unit at
$\mathfrak p_\theta$ with nonsquare residue. It cannot be a
square in $K$: a square which is a unit has a unit square root
and square residue. The Chinese remainder theorem imposes all
these congruences simultaneously and supplies the desired
arithmetic progression.

The construction is effective given exact number-field data:
enumerate rational primes, test complete splitting by finite
factorization outside the discriminant, compute residues, and
test squares in the finite residue fields. Chebotarev guarantees
that each required search terminates. No conjectural effective
prime bound is needed for this assertion.
\end{proof}

For the scalar formulas one may take $\Theta$ to consist of the
five rational contacts together with those roots of
$T^2+T+23$ which lie in $K$. Choose $C\ge2402$ in the
resulting progression. Then both $r_\beta$ and the normalizing
denominator are nonzero for every $X\in K$, and the conditional
definition over $\mathbb Q$ remains valid by
Proposition~\ref{p3:integer-shift}. The shift is allowed to depend
on $K$, but remains an ordinary fixed integer coefficient, not
an extra variable. The field-independent choice $2402$ suffices
when a real embedding exists.

No fixed integer shift can remove these contacts over every number
field. For any fixed $C$, the algebraic integer $X=\sqrt{-C}$ in
$\mathbb Q(\sqrt{-C})$ satisfies $X^2+C=0$. The field-dependent
choice in Lemma~\ref{p3:finite-square-avoidance} is therefore
essential if all number fields are to be considered.
